\scanpage{149}
% archive foliation: 148.
de $\mathrm{Sl}(2,\mathbb Z)$ sur $\mathcal D^+\times\mathbb C$. On peut en fait
faire opérer $\mathbb Z^2.\underbrace{\mathrm{Sl}(2,\mathbb Z).\deleted{\mathbb Z^2}}_{\text{\unsure{$\frac12$-direct}}}$
sur $\mathcal D^+\times\mathbb C$,
par
\[
T_{g,m,n}(\tau,z)=\left(g\tau,\frac z{c\tau+d}+m+ng\tau\right)
\tag{15}
\]
et l'opération \unsure{décrite|déduite} de $\mathrm{Sl}(2,\mathbb Z)$ sur
$\unsure{\mathcal X|\mathfrak X}_{\mathcal D^+}$ se déduit de (15) en passant au
quotient \deleted{\unsure{dans}} dans $\mathcal D^+\times\mathbb C$ par le sous-groupe
$\mathbb Z^2$ de $\mathbb Z^2.\mathrm{Sl}(2,\mathbb Z)$.

On peut énoncer dès maintenant

\uline{Théorème} La famille $\mathbb C$-analytique des courbes
elliptiques, \deleted{défini} \inserted{modulaire} pour la rigidification de Torelli
stricte, définie par
\diagram{
\begin{tikzcd}[ampersand replacement=\&]
\unsure{\mathcal X|\mathfrak X}_{\mathcal D^+}=\mathcal D^{\unsure{*|+}}\times\mathbb C/\mathbb Z^2 \arrow[d] \&
(\tau,z)\arrow[r,mapsto] \& \tau\\
\mathcal D^+ \&\&
\end{tikzcd}
}
\[
\mathbb Z^2\text{ opérant par}\qquad
T_{m,n}(\tau,z)=(\tau,z+m+n\tau)
\]
est un revêtement universel de $M_{1,1}$, avec
\deleted{\unsure{fibre de revêtement modulaire, et isomorphiques}}
comme groupe d'automorphismes sur $M_{1,1}$
le groupe $\mathrm{Sl}(2,\mathbb Z)$, opérant par
\[
g\tau=\frac{a\tau+b}{c\tau+d}
\qquad
g=\begin{pmatrix}a&b\\c&d\end{pmatrix},
\qquad
\begin{gathered}
a,b,c,d\in\mathbb Z\\
ad-bc=1.
\end{gathered}
\]
L'opération précédente est aussi décrite, en
termes du pb modulaire de Torelli, par
